%% LyX 2.3.6.2 created this file.  For more info, see http://www.lyx.org/.
%% Do not edit unless you really know what you are doing.
\documentclass[english]{article}
\usepackage[T1]{fontenc}
\usepackage[latin9]{inputenc}
\usepackage{babel}
\usepackage[unicode=true]
 {hyperref}
\begin{document}
\title{BiocPy: Porting Bioconductor\textquoteright s Data Representations
to Python }
\author{Jayaram Kancherla\textsuperscript{1}\\
\textsuperscript{1}Genentech Inc., South San Francisco, CA}
\maketitle
\begin{abstract}
Bioconductor is an open-source software project that provides tools
for the analysis and comprehension of genomic data. One of the main
advantages of Bioconductor is the availability of standard data representations
and large number of analysis tools for genomic experiments. These
tools allow researchers to efficiently store, manipulate, and analyze
their data, leading to a deeper understanding of the underlying biological
processes. However, Bioconductor data representations are currently
limited to the R programming language. BiocPy is an effort to port
these data representations in Python and integrate with existing scientific
ecosystem e.g. Numpy, Scipy etc. This includes core data structures
for managing genomic intervals (\emph{GenomicRanges}), container classes
for single (\emph{SummarizedExperiment}) or multi-assay experiments
(\emph{MultiAssayExperiment}), packages to efficiently perform operations
over matrices (\emph{mopsy}), interfaces to parse and represent genomic
experiments stored from R (\emph{rds2py}) and others. We also extend
\emph{SummarizedExperiment} to support single-cell experiments (\emph{SingleCellExperiment},
\emph{SpatialExperiment}). All packages in the BiocPy ecosystem are
published to Python's Package Index (PyPI). \href{http://github.com/BiocPy}{http://github.com/BiocPy}. 
\end{abstract}

\section{Introduction }

Bioconductor\cite{huber2015orchestrating} is an open-source community
effort to develop software packages for the analysis of genomic data.
It is widely used by researchers in the fields of genomics, bioinformatics,
and computational biology to perform statistical analyses, visualize
data, and develop predictive models. One of the key features of Bioconductor
is its core data structure and representations to efficiently manage
and represent various types of genomic experiments. However, these
data representations are currently limited to the R programming language,
which is a popular choice for statistical analysis but may not be
familiar or accessible to researchers in other domains. 

Python is a widely used programming language in the field of data
science, application or interface development, and has a large and
active community of users. Recent optimizations to the core of Python
has made the language 10-60\% faster for various workloads (\href{https://docs.python.org/3/whatsnew/3.11.html\#faster-cpython}{https://docs.python.org/3/whatsnew/3.11.html\#faster-cpython})
. Python has a number of libraries and tools that are specifically
designed for data analysis and has strong support for parallel computing,
which could be useful for working with large datasets. The machine
learning community has deeply invested in Python because of the availability
of libraries like keras, pytorch to rapidly prototype, build, publish
and deploy machine learning or artificial intelligence models. Researchers
often spend more time in the data gathering part of model development
because of the lack of an ecosystem of software packages to read,
parse and represent genomic experiments.

Inspired by Bioconductor, several isolated efforts implemented genomic
data representations in Python, e.g. the scverse ecosystem with AnnData
and MuData representations, pyRanges for representing genomic intervals
and genome annotations. However, these efforts are built for specific
technologies (single cell in the context of scverse) and may not provide
abstract classes to represent any type of genomic experiments. These
projects also lack compatibility with each other and provide fragmented
user experience across packages.

BiocPy is an effort to bridge the gap and port core data structures
and representations from Bioconductor in Python. tries to port familiar
well established and robust core set of data representations from
Bioconductor to Python. This includes core data structures for managing
genomic intervals and genome annotations (\emph{GenomicRanges}), container
classes for single (\emph{SummarizedExperiment}) or multi-assay experiments
(\emph{MultiAssayExperiment}), packages to efficiently perform operations
over matrices (\emph{mopsy}), interfaces to parse and represent genomic
experiments stored from R (\emph{rds2py}) and many others (\href{http://github.com/BiocPy}{http://github.com/BiocPy}.)
We also extend \emph{SummarizedExperiment} to support specialized
genomic experiments in single-cell (\emph{SingleCellExperiment}, \emph{SpatialExperiment}).
All packages in the BiocPy ecosystem are published to Python's Package
Index (PyPI). This effort to standardize representation of genomic
experiments can foster coordinated efforts across the community to
develop methods and tools for analysis and visualization of genomic
data.

\section{Core Data Representations in BiocPy }
\begin{itemize}
\item \textbf{Genomic Ranges:} One of the key data representations is the
\emph{GenomicRanges} package, which represents genomic intervals and
annotations in a compact and efficient manner. This package is particularly
useful for storing and manipulating large datasets, such as those
obtained from high-throughput sequencing experiments. It allows researchers
to easily select and subset specific intervals of interest, as well
as to perform various types of analyses on the data, such as overlaps
and intersections. 
\item \textbf{Summarized Experiments:} Another important data representation
is the \emph{SummarizedExperiment} package, which stores and summarizes
genomic data in a flexible and convenient manner. This package allows
researchers to store multiple types of data, including raw data, processed
data, and metadata, in a single object. It also provides tools for
data manipulation and visualization, making it easier for researchers
to explore and understand the experiment.
\end{itemize}

\section{Discussion }

\section{Software Availability }

\section{Acknowledgements}

I would like to thank Aaron Lun who has painstakingly implemented
rds2cpp, the C++ library to parse serialized S4 objects from RDS files.
I would also like to thank Hector Corrada Bravo, Michael Lawrence,Jean-Philippe
Fortin and William Hargreaves for many discussions related to various
packages.

\bibliographystyle{plain}
\bibliography{references}

\end{document}
